% Gerado a partir de MONOGRAFIA_PASE.md. O texto foi transposto sem alterações.

\chapter{Conclusão}\label{cap:6}

A IA deixou de ser apenas uma ferramenta de apoio e passou a atuar como agente capaz de planejar, decidir e executar trabalho de software com crescente independência. Este trabalho investigou o que essa mudança exige da arquitetura do software. A análise das cinco dimensões do ecossistema mostrou um padrão consistente: em todas elas, a capacidade técnica já demonstrada contrasta com a ausência de estruturas que a tornem confiável. Para responder a essa lacuna, o trabalho propôs a Arquitetura PASE, implementou uma prova de conceito e a aplicou à evolução de um produto real.

Quanto à primeira questão de pesquisa (QP1), a evidência analisada e a experiência da implementação convergem para a via híbrida. A delegação plena ainda é rara mesmo entre os usuários mais intensos de IA \cite{anthropic2026trends}, e a eficácia de sistemas multiagente depende da adequação entre a coordenação e a tarefa, e não do número de agentes \cite{kim2026scaling}. Além disso, os casos empresariais mais maduros combinam geração probabilística com execução determinística sob aprovações e auditoria \cite{vista2026outlook}. A própria validação reforçou essa conclusão: as falhas mais instrutivas da prova de conceito foram entregas plausíveis que só uma combinação de verificação por código e julgamento humano foi capaz de barrar. O cenário mais provável, e mais desejável, é o de ecossistemas em que a autonomia dos agentes é ampla e, ao mesmo tempo, subordinada a pontos de governança humana.

Quanto à segunda questão (QP2), a PASE responde com um desenho concreto, no qual supervisão, segurança e rastreabilidade não são requisitos não funcionais difusos, mas mecanismos estruturais distribuídos pelas três camadas. A prova de conceito demonstrou que esse desenho é viável. Uma demanda em linguagem natural percorreu as três camadas e chegou à produção com decisão humana em cada ponto crítico. A personalização em tempo de execução, por sua vez, adaptou a interface por usuário sem expor o produto a código arbitrário e recusou, no pior caso, conteúdos maliciosos que contornaram a primeira camada de validação.

\section{Contribuições}\label{sec:6.1}

A primeira contribuição deste trabalho é a análise multidimensional do software na era da IA agêntica, que reúne evidências de indústria, desenvolvimento, experiência do usuário, segurança e trabalho e as converte em cinco requisitos explícitos para arquiteturas de software agêntico. Essa conversão permite avaliar propostas futuras por critérios declarados, e não apenas por sua capacidade técnica.

A segunda contribuição é a própria Arquitetura PASE. Ela especifica, em três camadas, como demandas de origens heterogêneas são filtradas com memória do que já foi feito e recusado, e como a decisão humana é posicionada sobre a tarefa, o plano, a entrega e a publicação. Especifica também como a execução é isolada, verificada e revisada de forma independente, e como a personalização em tempo de execução é confinada a áreas com contratos verificáveis. Em relação aos trabalhos mais próximos, a arquitetura não apenas integra frentes antes tratadas isoladamente, mas as subordina a um mesmo regime de autonomia supervisionada.

A terceira contribuição é a prova de conceito e as evidências que ela produziu. A implementação conta com 533 testes automatizados, oito roteiros de demonstração executáveis e registros completos de cada decisão, da proposta bruta à publicação. Seus resultados documentam não apenas o caminho de sucesso, mas as falhas reais encontradas no caminho, como entregas aprovadas que não faziam o que foi pedido e estimativas de custo que erraram por mais de uma ordem de grandeza.

A quarta contribuição é um princípio de projeto que emergiu da validação e atravessa toda a arquitetura: garantias que dependem apenas de instrução ao modelo não são garantias. Em cada falha observada, a correção eficaz consistiu em transferir a garantia para uma verificação por código ou para uma decisão humana informada, com a regra de que o modelo aponta a evidência e o código confere se ela existe. Esse princípio é aplicável a outras arquiteturas que combinem componentes probabilísticos com responsabilidades críticas.

\section{Limitações}\label{sec:6.2}

As limitações do estudo começam pela revisão da literatura. Ela foi conduzida como revisão narrativa estruturada, e não como revisão sistemática com protocolo registrado, e parte do corpus é composta por versões preliminares e por relatórios de indústria, escolha justificada pela velocidade do fenômeno estudado, mas que reduz a robustez de algumas evidências. Vários relatórios de indústria, além disso, são produzidos por empresas que comercializam as tecnologias analisadas, e seus casos de uso provêm de clientes ou de empresas investidas. Esses números foram tratados como indicativos de tendência, e não como medições independentes.

A validação também tem alcance limitado. O produto-alvo é de porte reduzido, o número de tarefas executadas é pequeno e os operadores humanos foram os próprios autores, o que impede conclusões estatísticas sobre produtividade, sobre a carga cognitiva da supervisão ou sobre o comportamento da arquitetura em equipes reais. A personalização em tempo de execução foi validada quanto à segurança e ao funcionamento, mas não quanto ao seu efeito sobre a experiência de usuários reais. As métricas de supervisão, por sua vez, mostraram-se sensíveis à forma de cálculo, como evidenciou a diferença entre a taxa de intervenção na homologação e a taxa de aprovações revertidas.

Por fim, a prova de conceito implementa o núcleo da arquitetura, mas não sua totalidade. A prospecção competitiva, o barramento de eventos entre zonas, os fluxos declarativos de múltiplas etapas, as zonas conversacionais, os subgrafos de dependência, o replanejamento automático de tarefas dependentes e o agendamento periódico dos ciclos de manutenção permaneceram fora do escopo. O isolamento foi implementado por clone e ramo exclusivo, e não por contêiner efêmero por tarefa, e a integração com o GitHub, embora implementada, não foi exercitada na validação aqui relatada. Os modelos de linguagem usados evoluem rapidamente, de modo que resultados quantitativos de custo e de desempenho tendem a mudar com novas versões.

\section{Trabalhos Futuros}\label{sec:6.3}

O primeiro desdobramento é completar a implementação dos componentes que permaneceram fora do escopo, em especial o ambiente temporário de visualização para a homologação, a execução em contêineres efêmeros, a integração exercitada com repositórios remotos, as zonas sob demanda conversacionais e o barramento de eventos entre zonas. O ambiente de visualização é particularmente relevante, porque a validação mostrou que parte dos defeitos só se revela quando a pessoa observa o produto em execução.

O segundo desdobramento é uma avaliação em escala e em contexto real. Aplicar a arquitetura a produtos maiores, operados por equipes que não participaram de sua concepção, permitiria medir as métricas de supervisão com amostras significativas e comparar o desempenho com e sem os pontos de decisão humana. Permitiria também avaliar o efeito da personalização sobre usuários reais, inclusive quanto à transparência e à percepção de controle.

O terceiro desdobramento é a investigação dos efeitos sociotécnicos da arquitetura sobre as pessoas que a operam. A literatura indica que a preservação das competências depende de engajamento cognitivo ativo, e não da simples delegação à IA \cite{shen2026skill}. Investigar se o pacote de homologação da PASE, com critérios de aceite, evidências conferidas e histórico completo da interação com o modelo, favorece esse engajamento seria uma forma de avaliar se a arquitetura protege, além do produto, a capacidade humana de supervisioná-lo. Com isso, a PASE pretende oferecer um vocabulário e uma estrutura para que a confiabilidade seja tratada como requisito de projeto desde o início, e não como ajuste posterior.
